\textbf{T2: Thermal oscillations (10 pts)}

A resistor is made of a material which undergoes a phase transition so that
its resistance takes one of the two values, $R_1$ if its temperature is smaller
than $T_c$, and $R_2 > R_1$ if the temperature is larger than $T_c$.

\begin{center}\includegraphics[width=0.17\linewidth]{figures/EuPhO_2022_Q2_fig1.png}\end{center}

This resistor is connected to a voltage source through an inductor of
inductance $L$. It appears that if the applied voltage $V$ is between two
critical values, $V_1 < V < V_2$, the temperature of the resistor starts
oscillating. Assume that (i) the heat flux $P$ from the resistor to the ambient
medium is given by $P = \alpha(T - T_0)$, where $\alpha$ is a constant, $T$
denotes the temperature of the resistor, and $T_0$ is the ambient temperature;
(ii) the geometrical size of the resistor is so small that it will reach a
thermal equilibrium much faster than the characteristic time $L/R_2$.

\textbf{(a)} \textit{(2 pts)} Express $V_1$ and $V_2$ in terms of the other
parameters defined above.

\textbf{(b)} \textit{(6 pts)} Assuming that $V_1 < V < V_2$, sketch
qualitatively how the temperature of the resistor $T$ depends on time $t$, and
find the ratio $(T_\mathrm{max} - T_0)/(T_\mathrm{min} - T_0)$, where
$T_\mathrm{max}$ and $T_\mathrm{min}$ denote the maximal and minimal values of
$T$, respectively.

\textbf{(c)} \textit{(2 pts)} Find the period of oscillations if
$V = \sqrt{V_1 V_2}$ and $R_2 = 16R_1$.
